%epi_int_3d.tex
%
%mock ApJ paper on epi_int_3d: N-body simulations of spiral bending waves in planetary rings.
%Work in progress; only the vertical-motion section is drafted.
%
%AASTeX v7.0.2 (aastex702.cls, aasjournalv7.bst) is kept in this folder.
%to compile (tectonic runs bibtex itself):
%    tectonic epi_int_3d.tex
%jmh_bibliography.bib is kept local and not committed; cc_added.bib holds entries added by Claude Code.


\documentclass[twocolumn]{aastex702}

\usepackage{amsmath}

%journal
\received{not yet}
\revised{not yet}
\accepted{not yet}
\submitjournal{ApJ}

%short names
\shorttitle{Spiral Bending Waves in epi\_int\_3d}
\shortauthors{Hahn}


\begin{document}


%title page
\title{N-body Simulations of Spiral Bending Waves in Planetary Rings}

\author[gname=Joseph,sname=Hahn]{Joseph M. Hahn}
\affiliation{Space Science Institute}
\email[show]{jhahn@SpaceScience.org}
\correspondingauthor{Joseph M. Hahn}

\begin{abstract}
%placeholder
Draft. This mock paper currently holds only the derivation of the vertical
epicyclic motion used by the epi\_int\_3d drift step (Section \ref{sec:vertical}).
\end{abstract}


\section{Introduction}
\label{sec:intro}

%placeholder
The epi\_int and epi\_int\_lite N-body integrators \citep{HS13, H25} model a planetary ring as
a set of gravitating streamlines whose particles are confined to the ring plane. Spiral bending
waves, which are launched at a satellite's vertical resonances \citep{SCL83, S84}, require the
particles' vertical motions too. This paper describes epi\_int\_3d, a modest extension of
epi\_int\_lite into the vertical (out-of-plane) direction, and uses it to simulate spiral
bending waves.


\section{Vertical epicyclic motion}
\label{sec:vertical}

The epi\_int\_lite drift step advances each particle's unperturbed orbit about an oblate planet
using first-order epicyclic equations that are exactly reversible: converting spatial
coordinates into epicyclic orbit elements and back introduces no error, so the drift step
is not a source of numerical drift \citep{H25}. Here those equations are extended to
particles having vertical displacement $z$
and vertical velocity $v_z$, with the same first-order accuracy and the same exact
reversibility.

\subsection{Equations of motion}
\label{sec:eom}

A particle orbiting an oblate planet obeys Newton's second law,
\begin{equation}
    \label{eqn:newton}
    \ddot{\mathbf{r}} = -\nabla\Phi,
\end{equation}
where $\Phi$ is the planet's gravitational potential. That potential is a series in the
planet's zonal harmonics $J_{2k}$ (Eqn.\ 5 of \citealt{BL94}), and as in \cite{HS13} and
\cite{H25} only the $J_2$ term is kept, so
\begin{equation}
    \label{eqn:Phi3d}
    \Phi(r,z) = -\frac{GM_p}{\rho}\left[1 - J_2\left(\frac{R_p}{\rho}\right)^2P_2\left(\frac{z}{\rho}\right)\right].
\end{equation}
Here $M_p$ is the planet's mass, $R_p$ its radius, $r$ the particle's cylindrical radius,
$\rho^2 = r^2 + z^2$, and $P_2(\mu) = (3\mu^2 - 1)/2$. In cylindrical coordinates
$(r, \theta, z)$, Eqn.\ (\ref{eqn:newton}) is
\begin{eqnarray}
    \label{eqn:eom_cyl}
    \ddot{r} - r\dot{\theta}^2 &=& -\frac{\partial\Phi}{\partial r} \nonumber\\
    \frac{d}{dt}\left(r^2\dot{\theta}\right) &=& -\frac{\partial\Phi}{\partial\theta} = 0 \nonumber\\
    \ddot{z} &=& -\frac{\partial\Phi}{\partial z}.
\end{eqnarray}
The potential is independent of $\theta$, so the $z$-component of the specific angular
momentum, $h = r^2\dot{\theta} = rv_\theta$, is conserved exactly, even when $z \ne 0$.
The potential is also even in $z$, so near the ring plane
\begin{equation}
    \label{eqn:Phi_expand}
    \Phi(r,z) = \Phi(r,0) + \frac{1}{2}\nu^2(r)z^2 + {\cal O}(z^4),
\end{equation}
where $\Phi(r,0) = -GM_p/r - J_2GM_pR_p^2/2r^3$ is the midplane potential and
\begin{equation}
    \label{eqn:nu_r}
    \nu^2(r) = \left.\frac{\partial^2\Phi}{\partial z^2}\right|_{z=0}
        = \frac{GM_p}{r^3}\left[1 + \frac{9}{2}J_2\left(\frac{R_p}{r}\right)^2\right].
\end{equation}
Eqns.\ (\ref{eqn:eom_cyl}) then become
\begin{eqnarray}
    \label{eqn:eom_expand}
    \ddot{r} &=& \frac{h^2}{r^3} - \frac{\partial\Phi(r,0)}{\partial r} + {\cal O}(z^2) \nonumber\\
    \ddot{z} &=& -\nu^2(r)z + {\cal O}(z^3).
\end{eqnarray}

To first order the radial equation does not involve $z$; it is the planar equation of
motion. \cite{H25} solved the radial part of Eqns.\ (\ref{eqn:eom_expand}) to first order in
small displacements from circular motion, in a form that is exactly reversible: the
conversion from the coordinates $r, \theta, v_r, v_\theta$ to the epicyclic elements
$a, e, \tilde{\omega}, M$, and back, is exact. This is required if the solution is to be
used in the drift step of a drift-kick N-body integrator. That solution is summarized in
Appendix \ref{sec:appendix_inplane}. In it the particle's radius is $r = a[1 + {\cal O}(e)]$,
its mean longitude $\lambda = \tilde{\omega} + (\Omega_a/\kappa_a)M$ advances at the mean
angular velocity $\Omega_a$, and its mean anomaly $M$ advances at the epicyclic frequency
$\kappa_a$. Note that $\Omega$ denotes the longitude of ascending node, while $\Omega_a$ is
the mean angular velocity.

So the only additional physics that epi\_int\_lite's drift step needs in order to follow
out-of-plane motion is the vertical equation of motion, the second of
Eqns.\ (\ref{eqn:eom_expand}); the vertical forces that the ring and satellites exert enter
the kick step and are treated separately. Evaluating $\nu^2(r)$ at $r = a$ reduces the
vertical equation to the harmonic oscillator
\begin{equation}
    \label{eqn:z_oscillator}
    \ddot{z} = -\nu_a^2 z,
\end{equation}
with vertical frequency (Eqn.\ 17 of \citealt{BL94})
\begin{equation}
    \label{eqn:nu}
    \nu_a \equiv \nu(a) = \sqrt{\frac{GM_p}{a^3}}\left[1 + \frac{9}{2}J_2\left(\frac{R_p}{a}\right)^2\right]^{1/2},
\end{equation}
where the neglected terms are of order ${\cal O}(ei)$, due to the particle's radial excursion
$r - a = {\cal O}(ea)$, and ${\cal O}(i^3)$. Outside the planet the potential obeys Laplace's equation, which for an axisymmetric
potential is
\begin{equation}
    \label{eqn:laplace}
    \nabla^2\Phi = \frac{1}{r}\frac{\partial}{\partial r}\left(r\frac{\partial\Phi}{\partial r}\right)
        + \frac{\partial^2\Phi}{\partial z^2} = 0.
\end{equation}
At $r = a$ and $z = 0$, where $\Omega_a^2 = r^{-1}\partial\Phi/\partial r$,
$\kappa_a^2 = 3\Omega_a^2 + \partial^2\Phi/\partial r^2$ and
$\nu_a^2 = \partial^2\Phi/\partial z^2$, Eqn.\ (\ref{eqn:laplace}) relates the three
frequencies via $\kappa_a^2 + \nu_a^2 = 2\Omega_a^2$.

\subsection{Vertical solution}
\label{sec:elem2coords}

The solution of Eqn.\ (\ref{eqn:z_oscillator}) is a sinusoid of frequency $\nu_a$ whose
amplitude and phase define two new orbit elements, the inclination $i$ and the longitude of
ascending node $\Omega$:
\begin{eqnarray}
    \label{eqn:zvz}
    z   &=& a\,i\sin(\lambda - \Omega) \nonumber\\
    v_z &=& a\nu_a i\cos(\lambda - \Omega),
\end{eqnarray}
where the node regresses at the rate
\begin{equation}
    \label{eqn:node}
    \dot{\Omega} = \Omega_a - \nu_a \simeq -\frac{3}{2}J_2\left(\frac{R_p}{a}\right)^2\Omega_a.
\end{equation}
The vertical phase $\lambda - \Omega$ then advances at exactly $\nu_a$, so
Eqns.\ (\ref{eqn:zvz}) solve Eqn.\ (\ref{eqn:z_oscillator}) exactly, with $v_z = \dot{z}$,
and the particle crosses the ring plane upward at longitude $\lambda = \Omega$.
Measuring the phase from $\lambda$ rather than $\theta$ is what keeps the vertical motion
a pure harmonic oscillator; $\theta - \Omega$ would add an ${\cal O}(ei)$ error to the phase.

\subsection{Coordinates to elements}
\label{sec:coords2elem}

Eqns.\ (\ref{eqn:zvz}) are inverted by a fifth step that follows steps 1--4 of
Appendix \ref{sec:appendix_inplane}. Step 1 remains exact in three dimensions because $h$
is conserved and $r$ is the cylindrical radius, so steps 1--4 provide $a$, $\Omega_a$,
$\kappa_a$, $e$, $M$, $\tilde{\omega}$, and so $\lambda$. Step 5 computes $\nu_a$ from
Eqn.\ (\ref{eqn:nu}) and then
\begin{eqnarray}
    \label{eqn:i_node}
    i &=& \sqrt{\left(\frac{z}{a}\right)^2 + \left(\frac{v_z}{a\nu_a}\right)^2} \nonumber\\
    \Omega &=& \lambda - \text{atan2}(z, v_z/\nu_a).
\end{eqnarray}
Once $a$ is known, the mapping $(i, \lambda - \Omega) \leftrightarrow (z, v_z/\nu_a)$ is a
change between polar and Cartesian coordinates, so Eqns.\ (\ref{eqn:zvz}) and
(\ref{eqn:i_node}) are exact inverses. The drift step from $t$ to $t + \Delta t$ converts
coordinates to elements, advances $M \rightarrow M + \kappa_a\Delta t$ and
$\Omega \rightarrow \Omega + (\Omega_a - \nu_a)\Delta t$, and converts back. Its vertical part
is simply a rotation of $(z, v_z/\nu_a)$ through the angle $\nu_a\Delta t$.

\appendix

\section{The in-plane solution of Hahn et al.\ (2025)}
\label{sec:appendix_inplane}

This appendix summarizes the first-order, exactly reversible solution of the planar equation
of motion from Appendix A of \cite{H25}, in the notation used here. The mean angular velocity
and epicyclic frequency at semimajor axis $a$ are
\begin{eqnarray}
    \label{eqn:Omega_kappa}
    \Omega_a &=& \sqrt{\frac{GM_p}{a^3}}\left[1 + \frac{3}{2}J_2\left(\frac{R_p}{a}\right)^2\right]^{1/2} \nonumber\\
    \kappa_a &=& \sqrt{\frac{GM_p}{a^3}}\left[1 - \frac{3}{2}J_2\left(\frac{R_p}{a}\right)^2\right]^{1/2}.
\end{eqnarray}
In terms of the epicyclic elements $a$, $e$, $\tilde{\omega}$ and $M$, the particle's
coordinates are
\begin{eqnarray}
    \label{eqn:inplane}
    r &=& \frac{a}{\sqrt{1 + 2e\cos M}} \nonumber\\
    \theta &=& \frac{\Omega_a}{\kappa_a}(M + 2e\sin M) + \tilde{\omega} \nonumber\\
    v_r &=& ea\kappa_a\sin M\left(\frac{r}{a}\right)^3 \nonumber\\
    v_\theta &=& \frac{a^2\Omega_a}{r},
\end{eqnarray}
which conserve $h = rv_\theta = a^2\Omega_a$ exactly and solve the planar equation of motion
with errors of order ${\cal O}(e^2)$. The elements advance as $M \rightarrow M + \kappa_a\Delta t$,
with $a$, $e$ and $\tilde{\omega}$ constant. Eqns.\ (\ref{eqn:inplane}) are inverted exactly by:
\begin{enumerate}
\item Compute $h = rv_\theta$, then solve $h^2 = a^4\Omega_a^2$ for the semimajor axis,
    $a = R_p\left(C + \sqrt{C^2 - 3J_2/2}\right)$ where $C = h^2/2GM_pR_p$.
\item Compute $\Omega_a$ and $\kappa_a$ from Eqns.\ (\ref{eqn:Omega_kappa}).
\item Compute $e\sin M = (v_r/a\kappa_a)(a/r)^3$ and $e\cos M = \frac{1}{2}[(a/r)^2 - 1]$,
    which give $e$ and $M$.
\item Compute $\tilde{\omega} = \theta - (\Omega_a/\kappa_a)(M + 2e\sin M)$.
\end{enumerate}


\bibliography{jmh_bibliography,cc_added}{}
\bibliographystyle{aasjournalv7}  % v7.1 (June 2026) mislabels lone citations as "2006a"

\end{document}
